\section{Introduction}

Large language models are evaluated primarily on their ability to \emph{produce} working code: HumanEval~\cite{chen2021codex}, MBPP~\cite{austin2021mbpp}, LiveCodeBench~\cite{livecodebench2024}, and SWE-bench~\cite{jimenez2024swebench} all score whether a model can synthesize a function or patch that passes a test. This framing has driven impressive recent capability gains, but it leaves a substantively different question unanswered: \emph{can the model recognize, and structurally rewrite, an inefficiency in existing code that a production compiler cannot fix on its own?}

The question matters because the practical deployment surface for LLM-assisted optimization is precisely the gap between human-written and compiler-optimized code. A production C/C++ compiler at \texttt{-O2} already performs aggressive inlining, vectorization, peephole substitution, dead-store elimination, and (under \texttt{-flto}) cross-translation-unit motion. The marginal value of an LLM optimizer is only the speedup it captures \emph{beyond} that baseline. Yet every published benchmark we are aware of either omits the compiler baseline entirely or fixes the comparison at \texttt{-O0}, where the LLM's contribution is conflated with what the compiler would have done anyway.

We construct a benchmark whose \emph{selection criterion} is compiler-resistance. Every pattern in the suite is hand-validated against a five-regime sweep (\texttt{-O0}, \texttt{-O2 -fno-lto}, \texttt{-O3 -fno-lto}, \texttt{-O3 -fno-lto -ffast-math}, \texttt{-O3 -flto}) and against the LLVM polyhedral superoptimizer Polly. A variant is admitted only if the \texttt{fast.c} hand-optimized reference remains faster than \texttt{slow.c} by at least the pattern's expected lower-bound speedup after the compiler has done its best work. The compiler-resistance criterion is the discriminating axis the benchmark probes: 27 inefficiency patterns, each chosen because the optimization \emph{must} come from a semantic, algorithmic, or data-structural rewrite rather than from peephole substitution.

This paper makes the following contributions:

\begin{enumerate}
\item \textbf{A 1{,}244-variant compiler-resistant benchmark} (\S\ref{subsec:multiinput}, \S\ref{subsec:compiler-resistance}) spanning 27 base patterns across 7 categories ($391$ main variants), $675$ composed multi-pattern variants across $23$ two- and three-pattern combinations, and $178$ post-cutoff held-out variants ($36$ patterns at $5$ input-size / random-seed configurations each, minus $2$ small-$N$ sketch variants filtered as compiler-fixable). Of the $1{,}244$ active variants, $0$ are compiler-fixable under \texttt{-O3 -fno-lto} per the per-pattern threshold (\S\ref{subsec:compiler-resistance}); the $76$ variants that initially failed this gate ($49$ main, $25$ COMP, $2$ held-out) are preserved under \texttt{dataset/excluded/} with an \texttt{EXCLUDED.csv} manifest for transparency.
\item \textbf{A two-axis faithfulness verdict} (\S\ref{subsec:faithfulness}) that separates four model outcomes runtime correctness alone conflates: \textsc{Faithful} (the expected structural transformation), \textsc{Faithful-Alternative} (a different valid transformation that produces the same answer faster), \textsc{Structural-Only} (the structural shape was emitted but the candidate fails on adversarial inputs), and \textsc{Failed}. The verdict is computed by a cascade following \textsc{Invalidator}~\cite{lecong2023invalidator}: differential equivalence over nine input configurations followed by AST-based per-pattern structural matching with a regex fallback whose rate is reported as a data-quality signal.
\item \textbf{A setup-randomization protocol for defensible speedup claims} (\S\ref{subsec:stats}) following Mytkowicz et al.~\cite{mytkowicz2009producing}, which we apply to our own published numbers. Across six representative patterns spanning the speedup magnitude range, every pattern's $95\%$ bootstrap CI lower bound clears the $1.0$ gate, but CI \emph{width} varies from $5.3\%$ (algorithmically-dominated CF-3) to $51\%$ (microarchitecturally-dominated SR-1) --- a spread invisible to single-setup measurement.
\item \textbf{A held-out post-cutoff test set} (\S\ref{subsec:heldout}) of 178 variants across 36 patterns authored 2026-05-29 through 2026-05-31 and dated after every evaluated frontier model's training cutoff, following the temporal-boundary methodology of LiveBench~\cite{livebench2024} and SWE-bench-Live~\cite{swebenchlive2025}. Each pattern carries a primary-source citation (Redis HyperLogLog, zstd PRs, llama.cpp issues, rav1d, Fabian Giesen, Daniel Lemire, etc.) and a per-pattern expected speedup range honestly calibrated to the test hardware.
\item \textbf{A superoptimizer ablation} (\S\ref{subsec:compiler-resistance}) comparing the resistance of every variant against Polly's polyhedral pass in addition to standard compiler flags. Polly closes the slow/fast gap on $19\%$ of $1{,}176$ measured variants but is concentrated in two pattern families it was designed to handle (loop interchange for MI-4; loop fusion for HR-2, SR-3); the algorithmic (AL), input-sensitive (IS), data-structure (DS), human-style (HR), and control-flow (CF) categories remain largely Polly-resistant. Souper-class peephole superoptimizers are documented as out-of-scope by the Bansal--Aiken whole-SPEC ceiling of $1$--$5\%$~\cite{bansal2006gso}.
\item \textbf{A LoRA fine-tuning workflow with explicit held-out exclusion} (\S\ref{subsec:transfer}) so the held-out set never contaminates training data. Paired Wilcoxon signed-rank with $10{,}000$-resample bootstrap CI on the transfer-vs-baseline differential measures whether fine-tuning teaches a generalizing optimization skill or merely overfits to the training distribution.
\end{enumerate}

We adopt three design positions deliberately. First, we judge \emph{the artifact}, not the model's reasoning trace --- recent work~\cite{chen2025reasoningmodels, lanham2023measuring, turpin2023language} reports that frontier reasoning models verbalize the cues actually driving their answers only $25$--$39\%$ of the time on average, so conditioning the verdict on the model's stated plan would inherit that unreliability. Second, we report the LLM's contribution \emph{on top of} \texttt{-O2 -fno-lto}, the most common production default, rather than against \texttt{-O0}, so the speedup reflects real-world marginal value. Third, we treat any compiler-fixable variant as a benchmark bug and filter it (\texttt{scripts/filter\_non\_resistant.py}), preserving the variant for analysis under \texttt{dataset/excluded/} so future tooling can re-evaluate.

The remainder of the paper is organized as follows. Section~\ref{subsec:multiinput} describes the benchmark harness and multi-input robustness layer. Section~\ref{subsec:stats} details the statistical methodology and its application to our own results. Sections~\ref{subsec:faithfulness}--\ref{subsec:transfer} describe the faithfulness verifier, compiler-resistance validation, held-out construction, and cross-architecture infrastructure. We present results in Section~\ref{sec:results}, with reference to the related work in \texttt{related\_work.tex}.
